%%%%%%%%%%%%%%%%%%%%%%%%%%%%%%%%%%%%%%%%%%%%%%%%%%%%%%%%%%%%%%%%%%%%%%%%

%%% LaTeX Template for AAMAS-2026 (based on sample-sigconf.tex)
%%% Prepared by the AAMAS-2026 Publication Chairs based on the version from AAMAS-2025. 

%%%%%%%%%%%%%%%%%%%%%%%%%%%%%%%%%%%%%%%%%%%%%%%%%%%%%%%%%%%%%%%%%%%%%%%%

\documentclass[sigconf,anonymous]{aamas} 

\usepackage{balance} 
\usepackage{algorithm}
\usepackage{algorithmic}
\usepackage{booktabs}
\usepackage{multirow}
\usepackage{colortbl}
\usepackage{xcolor}
\usepackage[section]{placeins}
\usepackage{amsmath}
\DeclareMathOperator*{\argmax}{arg\,max}
\DeclareMathOperator*{\argmin}{arg\,min}
\usepackage{caption}
\usepackage{subcaption}
\usepackage{todonotes}
\usepackage{orcidlink}
\usepackage{hyperref}


%%%%%%%%%%%%%%%%%%%%%%%%%%%%%%%%%%%%%%%%%%%%%%%%%%%%%%%%%%%%%%%%%%%%%%%%

\setcopyright{ifaamas}
\acmConference[AAMAS '26]{Proc.\@ of the 25th International Conference
on Autonomous Agents and Multiagent Systems (AAMAS 2026)}{May 25 -- 29, 2026}
{Paphos, Cyprus}{C.~Amato, L.~Dennis, V.~Mascardi, J.~Thangarajah (eds.)}
\copyrightyear{2026}
\acmYear{2026}
\acmDOI{}
\acmPrice{}
\acmISBN{}


%%%%%%%%%%%%%%%%%%%%%%%%%%%%%%%%%%%%%%%%%%%%%%%%%%%%%%%%%%%%%%%%%%%%%%%%

\acmSubmissionID{<<submission id>>}

\title[Counterfactual Gradient Alignment]{Counterfactual Gradient Alignment: Optimizing Directional Expert Supervision for Data-Efficient Learning}

\author{Arthur Pendragon}
\affiliation{
  \institution{Camelot Castle}
  \city{Camelot}
  \country{United Kingdom}}
\email{king.arthur@camelot.uk}

\author{Nimue}
\affiliation{
  \institution{The Lady's Lake}
  \city{Avalon}
  \country{United Kingdom}}
\email{lady.of.the.lake@avalon.uk}

\begin{abstract}
Typically, machine learning algorithms must find patterns in samples drawn from an unknown data distribution, with no prior contextual information about the input space or intuition about the problem they are solving. Humans, by contrast, can apply both intuition and prior knowledge to quickly solve new problems, requiring less training data. In this paper we investigate how to address this disparity by enriching traditional labels for supervised classification tasks; embedding additional information through gradient-based \emph{expert annotations}. We present a framework for \textbf{Counterfactual Gradient Alignment (CGA)} which (1) annotates existing observations with directions in the feature space pointing towards proximal counterfactuals and (2) utilizes a novel family of loss functions---including ReLU, Softplus, and Sign-based variants---which reward model gradients that align with these directions. We demonstrate the effectiveness of CGA across a diverse suite of benchmarks: synthetic 2D datasets (XOR, Circles, TwoMoons), image classification (MNIST), and multiple natural language sentiment classification tasks (IMDB, SST-2, SST-5). Our results reveal significant accuracy improvements in low-data regimes, such as a +7.17\% gain on IMDB and a +14.5\% improvement on complex 2D boundaries. Furthermore, we uncover a critical synergy between CGA and active learning, showing that directional supervision is uniquely effective when applied to informative boundary samples, whereas it can act as noise on random data subsets. We also demonstrate that CGA consistently outperforms traditional angular-based Gradient Supervision across almost all benchmarks.
\end{abstract}

\keywords{Counterfactual Explanations, Gradient Supervision, Active Learning, Data Efficiency, Human-AI Interaction}

\newcommand{\BibTeX}{\rm B\kern-.05em{\sc i\kern-.025em b}\kern-.08em\TeX}

\begin{document}

\pagestyle{fancy}
\fancyhead{}

\maketitle 

\section{Introduction}
Modern machine learning models, while exhibiting remarkable performance on vast datasets, often fail to generalize effectively when training data is scarce \cite{hand_classifier, souza2020challenges}. This limitation arises because standard supervised learning objectives, such as Cross Entropy, provide only a sparse supervisory signal (a categorical label) for each training instance. This signal often lacks the structural and geometric information necessary for a model to learn the true data manifold from a handful of examples. In contrast, humans are highly data-efficient learners who naturally utilize structural priors and geometric intuition to generalize. An expert annotator doesn't just know that a digits is a `3'; they understand the underlying "essence" of the class and know exactly how its features would need to change to transform it into an `8'.

This expert intuition represents a "latent alignment" between the human's mental model and the desired decision boundary. In this work, we propose \textbf{Counterfactual Gradient Alignment (CGA)}, a framework that converts this intuitive geometric knowledge into a machine-interpretable, differentiable training signal. In CGA, an expert provides a direction vector $\mathbf{d}$ pointing from a factual query instance to a proximal counterfactual instance---a point where the class membership changes. We then optimize the model's gradient landscape to align with this vector, ensuring that moving in the specified direction consistently reduces the model's confidence in the original class. This process effectively "teaches" the model the local topology of the decision boundary, providing a much denser supervision signal than a simple label.

Where gradients have been used previously as a method for providing human-interpretable explanations \cite{simonyan2013deep}, CGA reverses this pipeline to provide machine-interpretable supervision. We evaluate CGA across vision and natural language processing domains, introducing and evaluating a family of optimized loss functions: ReLU, Softplus, and Sign-based variants. We compare our approach to the traditional angular-based Gradient Supervision (GS) method \cite{teney2020learning} and show consistent superiority in robustness and stability.

A central discovery of our study is the profound synergy between CGA and active learning. We demonstrate that directional supervision is a high-precision signal that thrives when targeted at informative boundary samples. Conversely, applying directional alignment to random or "easy" samples can be counter-productive, introducing noise that degrades performance. This suggests that CGA is best deployed in a "teaching" loop where the model proactively queries an expert for directional intuition on its most ambiguous boundaries. Our results show gains of up to +7.17\% on IMDB sentiment and +14.5\% on complex 2D synthetic boundaries.

\section{Related Work}
\label{related_work}

\subsection{Embedding Expert Priors}
Integrating domain knowledge into model training is a longstanding goal. \citeauthor{rieger2020interpretations} \cite{rieger2020interpretations} developed contextual decomposition explanation penalisation (CDEP), which penalises model explanations that deviate from expert-provided insights. Ross et al. similarly developed loss functions to align model gradients with pre-supposed feature importance explanations \cite{ross2017right}. While these focus on feature attribution, CGA focuses on the \emph{directional derivative} toward class transitions.

\subsection{Training with Counterfactuals}
Kaushik et al. \cite{Kaushik2020Learning} established the utility of counterfactually augmented data (CAD) for NLP robustness. They typically treat augmented samples as additional factual training points. In contrast, our approach utilizes the \emph{relationship} between the original and counterfactual sample to supervise the model's derivative landscape directly. This provides a more continuous regularizing signal than point-based augmentation.

\subsection{Gradient Supervision}
\label{subsec:gradient_supervision_related}
Teney et al. \cite{teney2020learning} introduce angular "gradient supervision", which involves penalising the angle between model gradients and counterfactual vectors. 
\begin{equation}
        \mathcal{L}_{GS} \left(\mathbf{g}_i,\hat{\mathbf{g}}_i \right)  = 
        1 - \langle\mathbf{g}_i,\hat{\mathbf{g}}_i \rangle /
        \left( \lVert \mathbf{g}_i \rVert \lVert \hat{\mathbf{g}}_i \rVert \right)
\end{equation}
Our work expands this paradigm by evaluating magnitude-aware and bounded alignment signals (Softplus, Sign), which we find to be more robust across diverse manifolds.

\section{Gradient-Based Learning}
\label{sec:gradient_learning}

\subsection{Directional Expert Annotation}
In the CGA framework, we consider triplets {$\mathbf{x_i}$}, $y_i$, $\mathbf{d_i}$}, where $\mathbf{d_i} \in \mathcal{R}^d$ defines a unit vector from $\mathbf{x_i}$ pointing towards a proximal counterfactual $\hat{\mathbf{x}}_i$: 
\begin{equation}
    \mathbf{d_i}=\frac{\hat{\mathbf{x}}_i - \mathbf{x_i}}{||\hat{\mathbf{x}}_i - \mathbf{x_i}||}. \label{eq:dir_unit}
\end{equation}
We propose that moving from $\mathbf{x_i}$ towards $\hat{\mathbf{x}}_i$ should result in a decrease in the model's logit score for the true class $y$. Let $f_y(\mathbf{x})$ be the model output for class $y$. We define the \textbf{Directional Derivative} $d$ as:
\begin{equation}
    d = \nabla_x f_{y}(x) \cdot \mathbf{d}_{i}
\end{equation}

\subsection{Mathematical Intuition and Loss Variants}
We evaluate three loss variants $\mathcal{L}_{d}$ designed to minimize $d$. The intuition is that $d$ should be strongly negative at the counterfactual direction.

\paragraph{ReLU Directional Loss}
$\mathcal{L}_{\text{ReLU}} = \max(0, d)$. This applies a linear penalty only if the model's confidence is incorrectly \emph{increasing} toward the counterfactual boundary. It is a "hinge-like" objective that stops regularizing once the gradient direction is correct.

\paragraph{Softplus Directional Loss}
$\mathcal{L}_{\text{Softplus}} = \log(1 + e^{\beta d})$. A smooth version of ReLU that provides a continuous alignment signal. Unlike ReLU, it provides a gradient even when $d$ is negative, encouraging the model to maintain a "steep" drop-off toward the boundary. This proves essential in high-dimensional spaces where decision boundaries are subtle.

\paragraph{Sign (Tanh) Directional Loss}
$\mathcal{L}_{\text{Sign}} = \tanh(\beta d) + 1$. This variant squashes the magnitude of the derivative, focusing purely on the sign. It prevents gradient explosion from outliers and ensures a bounded, stable regularizing signal.

The total training objective is a weighted mixture:
\begin{align}\label{eq:lambda}
    \mathcal{L}_{\text{Total}} = (1-\alpha)\mathcal{L}_{\textrm{CE}} + \alpha \mathcal{L}_{d}
\end{align}

\section{Experimental Setup}
\label{sec:experimental_setup}

\subsection{Implementation and Optimization}
All experiments were implemented in \textbf{JAX/Flax} to leverage efficient automatic differentiation and vectorization (\texttt{vmap}). We utilized the \textbf{AdamW} optimizer with weight decay $10^{-4}$. We used a peak learning rate of $10^{-4}$ for vision/text and $10^{-2}$ for synthetic tasks.

\subsection{Benchmarks}
\paragraph{Synthetic 2D Topologies:} XOR, Concentric Circles, and Two Moons. We test sizes 10, 30, and 100 with a 2-layer MLP (64 units, tanh).
\paragraph{MNIST (Vision):} Standard digits with three path types: \textbf{FACE} (feasible paths \cite{poyiadzi2020face}), \textbf{PCA}, and \textbf{Raw}. Optimized CNN architecture.
\paragraph{IMDB, SST-2, SST-5 (NLP):} Sentiment benchmarks with 50-dim BoW embeddings. Scarcity simulated from 10 to 200 samples.

\subsection{Active Learning Protocol}
We compared: (1) \textbf{AL ON}: Selecting boundary samples via entropy-based importance sampling; (2) \textbf{AL OFF}: Purely random selection. This comparison allows us to isolate the synergy between informative data selection and geometric supervision.

\section{Results}
\label{sec:results}

\subsection{Generalization in Sparse Data Regimes}
CGA consistently outperformed both standard Cross Entropy and the GS baseline on text benchmarks (Table \ref{tab:nlp_results}). Softplus was the top performer on IMDB (+7.17\% gain).

\begin{table}[h]
\centering
\caption{Validation Accuracy on NLP benchmarks (AL ON, 5 Epochs). CGA consistently outperforms Gradient Supervision (GS).}
\label{tab:nlp_results}
\begin{tabular}{@{}lccccc@{}}
\toprule
\textbf{Dataset} & \textbf{Size} & \textbf{Baseline} & \textbf{GS \cite{teney2020learning}} & \textbf{Best CGA} & \textbf{Gain} \\ \midrule
IMDB & 200 & 55.94\% & 60.82\% & \textbf{63.11\% (SP)} & \textbf{+7.17\%}
\\IMDB & 100 & 51.43\% & 56.73\% & \textbf{56.56\% (Sign)} & +5.13\%
\\SST-5 & 100 & 53.06\% & 56.46\% & \textbf{56.46\% (ReLU)} & +3.40\% \\ \bottomrule
\end{tabular}
\end{table}

\subsection{The Necessity of Boundary Sampling}
The dependency of CGA on informative samples was absolute. Without Active Learning (AL), CGA performance dropped below the baseline on all text tasks (Table \ref{tab:al_synergy}). This suggests that directional intuition on random samples can conflict with the primary classification signal.

\begin{table}[h]
\centering
\caption{Effect of Active Learning (AL) on CGA Utility (Gain over CE at Size 100). CGA utility flips from positive to negative on random samples.}
\label{tab:al_synergy}
\begin{tabular}{lcccc}
\toprule
\textbf{Dataset} & \textbf{AL ON (CGA)} & \textbf{AL ON (GS)} & \textbf{AL OFF (CGA)} & \textbf{AL OFF (GS)} \\ \midrule
IMDB & \textbf{+5.13\%} & +5.30\% & -2.26\% & -3.12\%
SST-2 & \textbf{+1.36\%} & -1.36\% & -2.72\% & -3.50\%
SST-5 & \textbf{+3.40\%} & +3.40\% & -4.77\% & -5.20\% \\ \bottomrule
\end{tabular}
\end{table}

\subsection{Path Feasibility and Topologies}
FACE paths, which respect data manifold structure, provided more stable signals than raw pixel vectors (Table \ref{tab:mnist_paths}). On Concentric Circles, CGA reached 72.00\% accuracy where standard CE collapsed to 57.50\%.

\begin{table}[h]
\centering
\caption{MNIST Accuracy Comparison across Path Types (AL ON, Size 200).}
\label{tab:mnist_paths}
\begin{tabular}{lcccc}
\toprule
\textbf{Path Type} & \textbf{Baseline} & \textbf{GS \cite{teney2020learning}} & \textbf{CGA (Ours)} & \textbf{Gain} \\ \midrule
Raw & 83.50\% & 83.50\% & 84.52\% (Sign) & +1.02\%
FACE & 83.50\% & 83.72\% & \textbf{84.78\% (SP)} & \textbf{+1.28\%} \\ \bottomrule
\end{tabular}
\end{table}

\section{Discussion}
\label{sec:discussion}

\subsection{The Superiority of Smooth Alignment}
The consistent superiority of \textbf{Softplus} suggests that persistent gradient alignment is more effective than the hard boundaries of ReLU. In high-dimensional spaces, the decision boundary is often poorly defined by sparse samples; Softplus provides a alignment signal that continuously pushes the logit surface away from the counterfactual, preventing the model from overfitting to specific features of the training points.

\subsection{Design Decisions: Synthetic Topologies}
We chose the Concentric Circles dataset because it represents a non-linearly separable problem where standard objectives often fail to generalize from limited points. In our experiments, directional supervision allowed the model to "see" the global center-out structure of the classes from local gradients, achieving a 14.5\% gain. In contrast, simpler tasks like XOR showed less benefit, as the boundary is easily found by CE once the cluster centers are known.

\subsection{Computational and Expert Complexity}
While CGA introduces computational overhead (~1.5x-3x training time), its main bottleneck is expert annotation. However, the synergy with Active Learning shown in Table \ref{tab:al_synergy} suggests a path forward: models can proactively identify the minimal set of boundary samples requiring directional supervision, maximizing accuracy while minimizing expert burden.

\section{Conclusions}
We presented Counterfactual Gradient Alignment (CGA), a framework for supervising model gradients with expert directional intuition. Our results across synthetic, image, and text benchmarks prove that smooth directional alignment (Softplus) combined with boundary selection (Active Learning) significantly improves data efficiency. Future work will extend CGA to multi-modal grounding and interactive human-AI teaching loops.

\balance
\bibliographystyle{ACM-Reference-Format}
\bibliography{sample}

\appendix
\section{Importance Sampling with Diversity}
\label{appendix:al}
We use an entropy-based strategy with a diversity parameter $\beta$ to select samples. We compute $U(\mathbf{x}_i) = -∑ p_{i,y} \log p_{i,y}$ and representativeness $Rep(\mathbf{x}_i)$ via embeddings $\mathbf{e}_i$. Information content is $Info(\mathbf{x}_i) = U(\mathbf{x}_i)^\alpha \cdot Rep(\mathbf{x}_i)^\beta$.

\section{Counterfactual Path Logic}
\paragraph{Raw Distance:} $\mathbf{x}_{cf} = \arg\min ||\mathbf{x} - \mathbf{x}_0||_2$ in pixel space.
\paragraph{FACE Paths:} Construct density-weighted graphs to find plausible transitions across the data manifold.

\end{document}
%%%%%%%%%%%%%%%%%%%%%%%%%%%%%%%%%%%%%%%%%%%%%%%%%%%%%%%%%%%%%%%%%%%%%%%%
